% संपूर्ण मराठी ग्रंथातील प्रकरण 62: प्रतिस्थापन
% हा संपादनयोग्य स्रोतखंड आहे. संकलनासाठी संपूर्ण स्रोतसंग्रहातील
% अक्षररूपे, पूर्वभाग, चिन्हव्याख्या आणि आकृती-स्रोत आवश्यक आहेत.
% मूळ: Open Logic Project; मजकूर आणि रूपांतर CC BY 4.0.
% यंत्रानुवाद, दुरुस्ती आणि तपासणी: OpenAI Codex —
% GPT-5.6 Sol आणि GPT-6 Sol, दोन्ही Ultra effort.
% दुरुस्ती-पुनरावलोकन: GPT-6.1 Sol, Ultra effort.
\chapter{प्रतिस्थापन}\label{sth:replacement:chap}




\section{परिचय}\label{sth:replacement:intro:sec}

प्रतिस्थापन हा स्वयंसिद्धक-रूपबंधच $\ZF$ आणि~$\Z$ यांतील फरक ठरवतो. \ref{sth:ordinals:chap}--\ref{sth:spine:chap} या प्रकरणांमध्ये आपण त्याचा वापर केला. आता अखेर या प्रकरणात आपण विचारू: प्रतिस्थापनाचे समर्थन करता येते का?

हा प्रश्न नेमका करण्यासाठी प्रतिस्थापन खरेच बरेच \emph{सामर्थ्यवान} आहे, हे लक्षात घ्यावे लागेल. ते किती सामर्थ्यवान आहे याची कल्पना या प्रकरणात (आणि पुन्हा \ref{sth:card-arithmetic:fix:sec} मध्ये) येईल. त्यावरून त्याच्या समर्थनाची खरोखर गरज आहे, हेही दिसेल.

समर्थनाचे दोन प्रकार पाहू. साधारणपणे, स्वयंसिद्धकाच्या फलितांवरून त्याचे समर्थन करण्याचा प्रयत्न म्हणजे \emph{बाह्य} समर्थन; तर संबंधित गणिती संकल्पनांतूनच स्वयंसिद्धकाला पाठबळ मिळते असे दाखवण्याचा प्रयत्न म्हणजे \emph{आंतरिक} समर्थन. याचा अर्थ या प्रकरणात अधिक स्पष्ट होईल; तरी ही चर्चा फार मोठ्या विषयाची केवळ सुरुवात आहे. अधिक माहितीसाठी विशेषतः मॅडी (1988a आणि 1988b) पाहा.







\section{प्रतिस्थापनाचे सामर्थ्य}\label{sth:replacement:strength:sec}

प्रतिस्थापनाच्या सामर्थ्याविषयी एक सोपे निरीक्षण आधी पाहू: $\Z$ च्या पुढे गेल्याशिवाय $\omega + \omega$ इतकी किंवा त्याहून मोठी कोणतीही फोन न्यूमन क्रमसूचक संख्या अस्तित्वात आहे, हे सिद्ध करता येत नाही.

का, याचा आराखडा असा. $\ZF$ मध्ये काम करताना $V_{\omega+\omega}$ हा संच पाहू. हा संच $\Z$ च्या एका \emph{प्रतिमानाचा} प्रांत म्हणून काम करतो. हे समजण्यासाठी सूत्राच्या \emph{सापेक्षीकरणाचा} संकेत ठरवू:
\begin{defn}\label{sth:replacement:strength:formularelativization}
	कोणताही संच $M$ आणि कोणतेही सूत्र $\phi$ दिल्यास, $\phi$ मधील सर्व संख्यापक $M$ पुरते मर्यादित केल्यावर मिळणारे सूत्र $\phi^M$ असे मानू. म्हणजे ``$\lexists[x]$'' ऐवजी ``$(\lexists[x \in M])$'' आणि ``$\lforall[x]$'' ऐवजी ``$(\lforall[x \in M])$'' लिहू.
\end{defn}\noindent
$\Z$ च्या प्रत्येक स्वयंसिद्धक $\phi$ साठी $\ZF \vdash
\phi^{V_{\omega+\omega}}$ हे दाखवता येते. पण \ref{sth:spine:rank:ordsetrankalpha} वरून $\omega+\omega$ ही संख्या $V_{\omega+\omega}$ मध्ये \emph{नसते}. म्हणून $\omega+\omega$ अस्तित्वात नसण्याशी $\Z$ सुसंगत आहे.

म्हणूनच \ref{sth:ordinals:replacement:sec} मध्ये आपण म्हटले की \ref{sth:ordinals:ordtype:thmOrdinalRepresentation} प्रतिस्थापनाशिवाय सिद्ध करता येत नाही. कारण $\Z$ मध्येच असा स्पष्ट सुक्रमण-संबंध परिभाषित करणे सोपे आहे की ज्याचा क्रमप्रकार सहजपणे $\omega+\omega$ \emph{असायला हवा}. नैसर्गिक संख्यांवरील पुढील क्रम दाखवताना \ref{sth:ordinals:idea:sec} मध्ये आपण याचे अनौपचारिक उदाहरण दिले होते:
\begin{align*}
	n \lessdot m \text{ तेव्हा आणि केवळ तेव्हाच }&\text{एकतर }n < m\text{ आणि }m-n\text{ सम आहे,}\\
	& \text{किंवा $n$ सम आणि $m$ विषम आहेत.}
\end{align*}
पण $\omega+\omega$ अस्तित्वात नसेल, तर हे सुक्रमण कोणत्याही क्रमसूचक संख्येशी समरूपी नसते. म्हणून $\Z$ वरून \ref{sth:ordinals:ordtype:thmOrdinalRepresentation} \emph{सिद्ध होत नाही}.

उलट पाहिल्यास: प्रतिस्थापनामुळे $\omega+\omega$ चे अस्तित्व, आणि म्हणून $V_{\omega+\omega}$ चे अस्तित्व, सिद्ध करता येते. एवढेच नव्हे. आपण परिभाषित करू शकणाऱ्या \emph{कोणत्याही} सुक्रमणासाठी \ref{sth:ordinals:ordtype:thmOrdinalRepresentation} सांगते की त्याच्याशी समरूपी अशी काही $\alpha$ असते; म्हणून $V_\alpha$ अस्तित्वात असते. अशा थेट अर्थाने प्रतिस्थापनामुळे संचांचा पदानुक्रम \emph{फार उंच} असण्याची हमी मिळते.

पुढील काही विभागांत आणि नंतर \ref{sth:card-arithmetic:fix:sec} मध्ये प्रतिस्थापनामुळे पदानुक्रम नेमका \emph{किती} उंच जातो याची अधिक चांगली कल्पना येईल. सध्या एवढाच साधा मुद्दा: प्रतिस्थापनाचे समर्थन \emph{खरोखर आवश्यक} आहे!







\section{प्रतिस्थापनाविषयी बाह्य विचार}\label{sth:replacement:extrinsic:sec}

आधी प्रतिस्थापनाचे \emph{बाह्य} समर्थन करण्याचा एक प्रयत्न पाहू. बूलोस यांनी असा युक्तिवाद सुचवला आहे:
\begin{quote}
  [\ldots] प्रतिस्थापनाची स्वयंसिद्धके स्वीकारण्याचे कारण अगदी सोपे आहे: त्यांचे अनेक इष्ट परिणाम आहेत आणि (दिसते तसे) अनिष्ट परिणाम नाहीत. क्रमिक संकल्पनेविषयीच्या प्रमेयांखेरीज त्या परिणामांत अनंत संख्यांची आदर्श नसली तरी समाधानकारक उपपत्ती आणि सुस्थापित संबंधांवरील विगमनाधारित व्याख्यांना समर्थन देणारा अत्यंत उपयुक्त निष्कर्ष येतो.
  (बूलोस 1971, पृ.~229)
\end{quote}
बूलोस यांच्या कल्पनेचा सारांश असा की प्रतिस्थापनाचे समर्थन त्याच्या फलितांवरून करावे. त्यांनी सांगितलेली विशिष्ट फलिते आपण मागील काही प्रकरणांत पाहिली आहेत. सुक्रमणाच्या क्रमप्रकारासाठी फोन न्यूमन क्रमसूचक संख्या उत्तम प्रतिनिधी ठरतात, हे प्रतिस्थापनामुळे सिद्ध करता आले (हीच आपली ``अनंत संख्यांची आदर्श नसली तरी समाधानकारक उपपत्ती''). प्रतिस्थापनामुळे $V_\alpha$ परिभाषित करता आले, प्रतांकाची संकल्पना मांडता आली आणि $\in$-विगमन सिद्ध करता आले (हीच ``क्रमिक संकल्पनेविषयीची प्रमेये''). अखेरीस, प्रतिस्थापनामुळे अनंतपार पुनरावर्तनाचे प्रमेय सिद्ध करता येते (हीच ``सुस्थापित संबंधांवरील विगमनाधारित व्याख्या'').

हे परिणाम खरोखर इष्ट आहेत. पण इतके इष्ट परिणाम प्रतिस्थापनाला \emph{समर्थन} देण्यासाठी पुरेसे आहेत का? \emph{नाही}. निदान सरळपणे तरी नाही.

एक सोपी अडचण अशी आहे: प्रतिस्थापनाचे काही इष्ट परिणाम आपण मांडले, पण त्यांपैकी अनेक परिणाम दुसऱ्या मार्गांनीही मिळू शकले असते. ही बाब जितकी माहीत असायला हवी तितकी नाही; म्हणून परिस्थिती स्पष्ट करू.

संचांची एक सोपी उपपत्ती आहे: स्तर उपपत्ती, संक्षेपाने $\LT$.\footnote{$\LT$ ची पहिली रूपे मॉन्टेग्यू (1965) आणि स्कॉट (1974) यांनी दिली. पॉटर (2004) यांनी ती सोपी करून पुस्तकरूपात सविस्तर मांडली; अलीकडे बटन (2021) यांनी $\LT$ आणखी सोपी केली.} $\LT$ ची स्वयंसिद्धके फक्त विस्तारात्मकता, विभक्तीकरण आणि प्रत्येक संच कोणत्या तरी \emph{स्तराचा} उपसंच असतो हे विधान इतकीच आहेत. येथे ``स्तर'' अशी कल्पक व्याख्या आहे की हे स्तर आपल्या परिचयाच्या $V_\alpha$ प्रमाणे वागतात. म्हणून $\ZF$ वरून $\LT$ सिद्ध होते; पण $\LT$ ही $\ZF$ पेक्षा \emph{खूप} दुर्बल आहे. खरे तर $\LT$ वरून जोड्या, घातसंच, अनंतता किंवा प्रतिस्थापन यांपैकी काहीही मिळत नाही. $\LT$ मध्ये अनंतता आणि घातसंच जोडून मिळणाऱ्या उपपत्तीला $\Zr$ म्हणू. यावरून जोड्याही मिळतात; म्हणून $\Zr$ निदान $\Z$ इतकी सामर्थ्यवान आहे. पण प्रत्यक्षात $\Zr$ ही $\Z$ पेक्षा काटेकोरपणे अधिक सामर्थ्यवान आहे; कारण प्रत्येक संचाला प्रतांक असतो हे विधान ती जोडते (म्हणूनच तिला $\Zr$ म्हणावे असे मी सुचवतो). $\Zr$ वरून क्रमसूचक संख्यांची पूर्णपणे समाधानकारक उपपत्ती, पदानुक्रमाचे सुक्रमित टप्प्यांत विभाजन करणारे निष्कर्ष, $\in$-विगमनाचा पुरावा आणि अनंतपार पुनरावर्तनाचे एक \emph{रूप} मिळते.

थोडक्यात, बूलोस यांना ही बाब माहीत नसली, तरी त्यांनी सांगितलेले सर्व इष्ट परिणाम प्रतिस्थापनाशिवाय \emph{मिळू शकले असते}; त्यासाठी $\Z$ ऐवजी $\Zr$ वापरणे पुरेसे होते.

(मग आपण प्रचलित पद्धत वापरून $\LT$ आणि $\Zr$ ऐवजी $\ZF$ का शिकवली? दोन कारणे आहेत. पहिले: निव्वळ ऐतिहासिक कारणांमुळे $\LT$ पासून सुरुवात करणे अपारंपरिक आहे; संच उपपत्तीवरील अधिक प्रमाणित चर्चा वाचता याव्यात अशी आपली इच्छा होती. दुसरे: $\LT$ आणि $\Zr$ समजून घेण्याची तयारी झाल्यावर पॉटर (2004) आणि बटन (2021) थेट वाचता येतील.)

अर्थात $\Zr$ ही $\ZF$ पेक्षा काटेकोरपणे दुर्बल असल्याने, $\ZF$ सिद्ध करते पण $\Zr$ अनिर्णित ठेवते असे निष्कर्ष आहेत. त्यांपैकी एखाद्या निष्कर्षाच्या आधारे प्रतिस्थापनाचे बाह्य समर्थन करण्याचा प्रयत्न करता येईल. पण $\Zr$ कशी वापरावी हे समजल्यावर (अ) प्रतिस्थापनामुळेच, दुसऱ्या मार्गाने नाही, ठरणाऱ्या आणि (ब) सहज अर्थाने खऱ्या वाटणाऱ्या विधानांची अनेक उदाहरणे शोधणे बरेच कठीण असते. (अधिक माहितीसाठी पॉटर (2004, \S13.2) पाहा.)

मुख्य निष्कर्ष असा: प्रतिस्थापनाचे पटणारे बाह्य समर्थन द्यायचे असेल, तर प्रतिस्थापनाशिवाय \emph{मिळूच शकणार नाही} असा निष्कर्ष शोधावा लागेल. हे सोपे काम नाही.

प्रतिस्थापनाच्या निव्वळ बाह्य समर्थनाच्या कोणत्याही प्रयत्नासमोर आणखी एक अडचण उभी राहते; कदाचित ती पहिल्यापेक्षा अधिक मूलभूत आहे. बूलोस फक्त प्रतिस्थापनाचे अनेक इष्ट परिणाम सांगत नाहीत; त्याचे ``(दिसते तसे) अनिष्ट'' परिणाम नाहीत, असेही म्हणतात. पण कंसातील ``दिसते तसे'' ही सावधगिरीची नोंद अतिशय महत्त्वाची आहे.

आपण येथे कसे पोहोचलो ते आठवू. अनौपचारिक गुणधर्माधारित संचरचनेत विसंगती निर्माण झाली; त्याला प्रतिसाद म्हणून आपण संचांची संचयी-क्रमिक संकल्पना स्वीकारली. या संकल्पनेबरोबर एक मांडणी येते आणि ती आपल्याला तिच्या सुसंगततेची खात्री देईल, अशी आशा आहे. पण त्या मांडणीच्या आतून प्रतिस्थापनाचे समर्थन करता आले नाही, तर $\ZF$ सुसंगत आहे असे मानण्याचे (अजून) कोणतेही कारण आपल्याकडे नाही. किंवा अधिक नेमके: आत्तापर्यंत कोणालाही \emph{विसंगती सापडलेली नाही} या (कदाचित निव्वळ योगायोगाच्या) वस्तुस्थितीखेरीज $\ZF$ च्या सुसंगततेवर विश्वास ठेवण्याचे कारण नाही. नेमक्या याच अर्थाने बूलोस यांचे विधान ``(दिसते तसे) $\ZF$ सुसंगत आहे'' इतकेच ठरते. सुसंगततेबाबत याहून अधिक खात्री आपण मागितली पाहिजे.

$\ZF$ च्या (ज्ञात) परिणामांपुरतेच मर्यादित असलेले, म्हणजे प्रतिस्थापनाचे \emph{निव्वळ} बाह्य समर्थन करण्याचे, कोणतेही प्रयत्न या अडचणीला सामोरे जातील. म्हणून आपण प्रतिस्थापनाचे \emph{आंतरिक} समर्थन शोधू: संचांविषयी आपण सांगितलेली मांडणी आपल्याला आधीपासूनच प्रतिस्थापन स्वीकारण्यास बांधील करते, असे सुचवणारे समर्थन.








\section{आकारमानाच्या मर्यादेचे तत्त्व}\label{sth:replacement:limofsize:sec}

प्रतिस्थापनाचे ``आंतरिक'' समर्थन देण्याचा बहुधा सर्वाधिक प्रचलित प्रयत्न पुढील कल्पनेवर आधारलेला असतो:
\begin{enumerate}
  \item[] \limofsize. कोणत्याही वस्तूंचा संच तयार होतो, जर त्या फार जास्त नसतील तर.
\end{enumerate}
आकारमानाच्या पुढील तुलनेसाठी निवडीचे स्वयंसिद्धक गृहीत धरून, हे तत्त्व प्रतिस्थापनाला लगेच समर्थन देईल. कारण प्रतिस्थापनाने तयार झालेला कोणताही संच ज्या संचापासून तयार झाला त्याच्यापेक्षा मोठा असू शकत नाही. नेमके सांगायचे तर: समजा प्रतिस्थापन वापरून तुम्ही
$\funimage{\tau}{A} = \Setabs{\tau(x)}{x \in A}$ हा संच तयार केला; मग $\cardle{\funimage{\tau}{A}}{A}$; म्हणून $A$ चे सदस्य संच तयार होण्याइतके कमी असतील, तर त्यांच्या प्रतिमाही $\funimage{\tau}{A}$ तयार होण्याइतक्या कमी असतील.

प्रतिस्थापनाचे समर्थन करण्यासाठी \limofsize{} वापरण्यातील उघड अडचण अशी की \limofsize{} सारखे कोणतेही तत्त्व आपण अजून मांडलेलेच \emph{नाही}. शिवाय, \ref{sth:story:chap}--\ref{sth:z:chap} मध्ये संचांच्या संचयी-क्रमिक संकल्पनेची मांडणी करताना \limofsize{} कडे निर्देश करणारा \emph{संकेतही} त्यात नव्हता. म्हणूनच बूलोस यांनी एकदा असे लिहिले: ``संच उपपत्तीच्या `मागे' निदान दोन विचार आहेत, असा निष्कर्ष कदाचित काढता येईल'' (1989, पृ.~19). एका बाजूला, संचांच्या संचयी-क्रमिक संकल्पनेशी संबंधित कल्पना $\Z$ ला समर्थन देण्यासाठी आहेत. दुसऱ्या बाजूला, \limofsize{} हे प्रतिस्थापनाला समर्थन देण्यासाठी आहे.

पण अडचण एवढीच नाही की \limofsize{} बाबत आपण आतापर्यंत \emph{मौन} बाळगले आहे. उलट, नुकतेच मांडलेले \limofsize{} तत्त्व संचांच्या संचयी-क्रमिक संकल्पनेशी जुळणेही अवघड वाटते. कारण संच \emph{टप्प्यांत} तयार होतात, या कल्पनेचा त्यात अजिबात उल्लेख नाही.

या ``दोन विचारांचा'' (म्हणजे संचांची संचयी-क्रमिक संकल्पना आणि \limofsize) इतिहास पाहता हे फारसे आश्चर्यकारक नाही. संच उपपत्तीतील विरोधाभास रोखण्याचे हे मुळात दोन भिन्न प्रकल्प आहेत. संचांची संचयी-क्रमिक संकल्पना रसेलचा विरोधाभास साधारणपणे असे म्हणून रोखते: \emph{रसेलचा संच अस्तित्वात असेल अशी अपेक्षाच आपण करायला नको होती, कारण तो कोणत्याही टप्प्यावर ``तयार'' होणार नाही}. याउलट, \limofsize{} रसेलचा संच साधारणपणे असे म्हणून नाकारते: \emph{रसेलचा संच अस्तित्वात असेल अशी अपेक्षाच आपण करायला नको होती, कारण तो फार मोठा असता}.

असे मांडल्यावर स्पष्टच सांगू: विरोधाभासांना उत्तर म्हणून \limofsize{} ला \emph{खूपच} अधिक समर्थन आवश्यक आहे. उदाहरणार्थ, रसेलच्या विरोधाभासाचे हे रूप पाहा: \emph{स्वतःचाच वास न घेणाऱ्या पग कुत्र्यांचाच नेमका वास घेणारा कोणताही पग कुत्रा नाही} (\ref{sth:story:rus:sec} पाहा). ``असा पग कुत्रा का नाही?'' असे विचारल्यावर त्याला फारच जास्त पग कुत्र्यांचा वास घ्यावा लागला असता, हे चांगले उत्तर ठरणार नाही. मग रसेलचा संच अस्तित्वात नाही याचे सहज पटणारे स्पष्टीकरण तो अस्तित्वात असण्यासाठी ``फार मोठा'' झाला असता, हे कसे ठरेल?

थोडक्यात, \limofsize{} मागील ``सहजस्फूर्त'' प्रेरणेविषयी तुम्हाला थोडे गूढ वाटले, तर ते समजण्यासारखे आहे.







\section{प्रतिस्थापन आणि ``निरपेक्ष अनंतता''}\label{sth:replacement:absinf:sec}

आता \limofsize{} बाजूला ठेवून प्रतिस्थापनाचे समर्थन करण्याचे वेगळ्या प्रकारचे (आंतरिक) प्रयत्न पाहू. हे प्रयत्न संच टप्प्याटप्प्याने तयार होतात या कल्पनेला गांभीर्याने घेतात.

क्रमिक प्रक्रिया प्रथम सांगताना आपण प्रत्येक टप्प्यावर काय घडते हे स्पष्ट करणारी काही तत्त्वे मांडली होती: \stageshier, \stagesord आणि \stagesacc. नंतर टप्प्यांच्या संख्येविषयी काही सांगणारी तत्त्वे जोडली: \stagessucc{} नुसार संच तयार होण्याची प्रक्रिया कधीही संपत नाही, आणि \stagesinf{} नुसार ती अनंताव्या टप्प्यापर्यंत पोहोचते.

ही नंतरची दोन तत्त्वे आणखी व्यापक तत्त्वातून निष्पन्न होतात असे सुचवणे रास्त आहे; उदाहरणार्थ:
\begin{enumerate}
  \item[] \stagesinex. टप्प्यांची संख्या निरपेक्ष अर्थाने अनंत आहे; पदानुक्रम जितका उंच असू शकतो तितका उंच आहे.
\end{enumerate}
अर्थात हे एक अनौपचारिक तत्त्व आहे. पण संचांच्या संचयी-क्रमिक संकल्पनेतून ते लगेच \emph{निष्पन्न} होत नसले, तरी त्या संकल्पनेशी निश्चितच \emph{सुसंगत} वाटते. निदान \limofsize{} च्या विपरीत, संच टप्प्याटप्प्याने तयार होतात ही कल्पना त्यात टिकून राहते.

आता \stagesinex{} वापरून प्रतिस्थापनाचे समर्थन करता येईल अशी आशा आहे. ते कसे शक्य होईल ते पाहू.

\ref{sth:ordinals:idea:sec} मध्ये आपण पाहिले की (अनौपचारिक अर्थाने) $\omega+\omega$ शी समरूप असायला हवे असे सुक्रमण सहज तयार करता येते. दुसऱ्या शब्दांत, (अनौपचारिक अर्थाने) $\omega+\omega$ हा क्रमप्रकार असलेली टप्प्याटप्प्यांची क्रमिक प्रक्रिया आपण सहज कल्पू शकतो. म्हणून \stagesinex{} स्वीकारले असेल, तर पदानुक्रमात किमान $\omega+\omega$-वा टप्पा, म्हणजे $V_{\omega+\omega}$, आहे हेही स्वीकारले पाहिजे; कारण पदानुक्रम इतक्यापर्यंत \emph{नक्कीच जाऊ शकतो}.

हा विचार सर्वसाधारणपणे असा मांडता येतो: कोणतेही सुक्रमण घ्या; क्रमिक पदानुक्रम बांधण्याची प्रक्रिया किमान त्या सुक्रमणाइतकी पुढे गेली पाहिजे. \ref{sth:ordinals:ordtype:thmOrdinalRepresentation} याला \emph{स्वयंसिद्धक} मानल्यास याची खात्री मिळेल. मग प्रत्येक सुक्रमण एखाद्या फोन न्यूमन क्रमसूचक संख्येशी समरूप आहे असे ठरेल. प्रत्येक फोन न्यूमन क्रमसूचक संख्येचा प्रतांक ती स्वतःच असल्याने, \ref{sth:ordinals:ordtype:thmOrdinalRepresentation} यावरून असे ठरेल की संच उपपत्तीत आपण कोणतेही सुक्रमण वर्णन करू शकलो, तर संच बांधण्याची क्रमिक प्रक्रिया त्या सुक्रमणाच्याही पुढे गेली पाहिजे.

हा विचार \stagesinex{} चा परिणाम वाटतो. दुर्दैवाने, त्यावरून प्रतिस्थापन मिळवायचे असेल, तर एक सोपा तांत्रिक अडथळा आहे: प्रतिस्थापन हे \ref{sth:ordinals:ordtype:thmOrdinalRepresentation} पेक्षा काटेकोरपणे अधिक सामर्थ्यवान आहे. (ही बाब पॉटर (2004, \S13.2) यांनी नोंदवली आहे; तिचा पुरावा आपण \ref{sth:replacement:finiteaxiomatizability:sec} मध्ये देऊ.)

म्हणून \stagesinex{} चा अर्थ प्रतिस्थापन मिळेल अशा रीतीने घ्यायचा असेल, तर पदानुक्रम प्रत्येक सुक्रमणाच्या पुढे जातो \emph{एवढेच} त्याने सांगून चालणार नाही. त्याहून बलवान दावा आवश्यक आहे. यासाठी शोनफेल्ड (1977) यांनी या कल्पनेचा एक अत्यंत स्वाभाविक विस्तार सुचवला: पदानुक्रम कोणत्याही संचाशी \emph{सहअंतिम} नाही.\footnote{G\"odel यांनीही असाच विचार सुचवला असावा; पॉटर (2004, पृ.~223) पाहा. G\"odel आणि शोनफेल्ड यांच्या चर्चेसाठी इनकुर्वाती (2020, पृ.~90--95) पाहा.} अधिक स्पष्टपणे: $\tau$ हे संचांना पदानुक्रमातील टप्प्यांकडे पाठवणारे प्रतिचित्रण असेल, तर कोणत्याही संच $A$ ची $\tau$ खालील प्रतिमा संपूर्ण पदानुक्रम संपवू शकत नाही. दुसऱ्या शब्दांत (रूपबंधाच्या स्वरूपात):
\begin{enumerate}
  \item[] \stagescofin. $A$ हा संच असेल आणि प्रत्येक $x \in A$ साठी $\tau(x)$ हा टप्पा असेल, तर प्रत्येक $x \in A$ च्या $\tau(x)$ नंतर येणारा एक टप्पा असतो.
\end{enumerate}
$\ZF$ मध्ये \stagescofin{} चे योग्य औपचारिक रूप सिद्ध होते, हे स्पष्ट आहे. उलट, \stagescofin{} प्रतिस्थापनाचे समर्थन करते, असा अनौपचारिक युक्तिवाद करता येतो.\footnote{\stagescofin{} च्या एखाद्या औपचारिक रूपावरून प्रतिस्थापन सिद्ध करणे अधिक कठीण ठरेल; कारण केवळ $\Z$ मध्ये टप्प्यांची व्याख्या करण्याइतके सामर्थ्य नाही. त्यामुळे \stagescofin{} चे औपचारिकीकरण कसे करावे हे स्पष्ट नाही. मात्र \ref{sth:replacement:extrinsic:sec} मध्ये चर्चिल्याप्रमाणे $\LT$ च्या एखाद्या विस्तारात काम करणे हा एक पर्याय आहे.} समजा $(\forall x \in A)\lexists![y][\phi(x,y)]$. मग प्रत्येक $x \in A$ साठी $\phi(x,y)$ पूर्ण करणारा $y$ म्हणजे $\sigma(x)$ मानू, आणि $\sigma(x)$ पहिल्यांदा ज्या टप्प्यावर तयार होतो तो $\tau(x)$ मानू. \stagescofin{} नुसार असा टप्पा $V$ आहे की $(\forall x \in A)\tau(x)\in V$. आता प्रत्येक $\tau(x) \in V$ आणि $\sigma(x) \subseteq \tau(x)$ असल्याने, विभक्तीकरण वापरून $\Setabs{y
\in V}{(\exists x \in A)\sigma(x) = y} = \Setabs{y}{(\exists x \in
A)\phi(x,y)}$ हा संच मिळवता येतो.

\begin{prob}
  $\ZF$ मध्ये \stagescofin{} चे औपचारिकीकरण करा.
\end{prob}

म्हणून \stagescofin{} प्रतिस्थापनाला समर्थन देते. आणि \stagesinex{} हे \stagescofin{} ला समर्थन देते, हे निदान संभवनीय वाटते. समजा \stagescofin{} खोटे ठरते. म्हणजे पदानुक्रम एखाद्या संच~$A$ शी सहअंतिम आहे: असे $\tau$ प्रतिचित्रण आहे की प्रत्येक टप्पा~$S$ साठी असा काही $x \in A$ आहे की $S \in \tau(x)$. अशा परिस्थितीत पदानुक्रमाची कथित ``निरपेक्ष अनंतता'' पकडण्याचा मार्ग आपल्याकडे आहे: $A$ वर लागू केलेल्या $\tau$ च्या परासाने ती \emph{संपूर्णपणे व्यापली} जाते. त्यामुळे पदानुक्रम ``निरपेक्षपणे अनंत'' आहे या विचाराला बाधा येते. व्यत्यासाने: \stagesinex{} वरून \stagescofin{} मिळते; आणि ते प्रतिस्थापनाचे समर्थन करते.

प्रतिस्थापनाचे आंतरिक समर्थन देण्याचा हा खरोखर आशादायक प्रयत्न आहे. मात्र तो अखेरीस सफल होतो की नाही, हे ठरवण्याचे काम आम्ही तुम्हालाच सोपवतो.







\section{प्रतिस्थापन आणि प्रतिबिंबन}\label{sth:replacement:ref:sec}

\stagesinex{} च्या साहाय्याने प्रतिस्थापनाचे समर्थन करण्याचा आपला शेवटचा प्रयत्न एका सखोल आणि सुंदर निष्कर्षापासून सुरू होतो:\footnote{आठवण: स्पष्टपणे अन्यथा म्हटले नसेल, तर सर्व सूत्रांमध्ये प्राचले असू शकतात.}


\begin{thm}[प्रतिबिंबन रूपबंध]\label{sth:replacement:ref:reflectionschema}
कोणत्याही सूत्र $\phi$ साठी:
\[
\forall \alpha \exists \beta > \alpha (\forall x_1 \ldots, x_n \in
V_\beta)(\phi(x_1, \ldots, x_n) \liff \phi^{V_\beta}(x_1, \ldots, x_n))
\]
\end{thm}
\noindent
\ref{sth:replacement:strength:formularelativization} प्रमाणे, $\phi$ मधील प्रत्येक संख्यक $V_\beta$ या संचापुरता मर्यादित केल्यावर $\phi^{V_\beta}$ मिळते. म्हणून सहज अर्थाने प्रतिबिंबन असे सांगते: $\phi$ संपूर्ण पदानुक्रमात सत्य असेल, तर पदानुक्रमाच्या हवे तितक्या अनेक \emph{आरंभीच्या खंडांमध्ये}ही $\phi$ सत्य असते.

मॉन्टेग्यू (1961) आणि लेव्ही (1960) यांनी दाखवले की (योग्यरीत्या मांडलेली) प्रतिस्थापन आणि प्रतिबिंबन ही $\Z$ च्या साहाय्याने समतुल्य आहेत; म्हणजे त्यांपैकी कोणतेही एक जोडल्यास $\ZF$ मिळते. (हे निष्कर्ष आपण \ref{sth:replacement:refproofs:sec} मध्ये सिद्ध करू.) ही समतुल्यता लक्षात घेऊन, \stagesinex{} वरून प्रतिबिंबन आणि प्रतिस्थापनाचे समर्थन करण्याचा पुढील प्रयत्न करता येईल: \stagesinex{} असेल, तर पदानुक्रम खूपच उंच असायला हवा; इतका उंच की त्याविषयी आपण म्हणू शकतो अशा कोणत्याही गोष्टीने त्याच्या उंचीला बंध घालता कामा नये. हा विचार आपण असा समजू शकतो: कोणतेही वाक्य~$\phi$ संपूर्ण पदानुक्रमात सत्य असेल, तर ते पदानुक्रमाच्या हवे तितक्या अनेक आरंभीच्या खंडांतही सत्य असते. हेच प्रतिबिंबन आहे.

पुन्हा, प्रतिस्थापनाचे आंतरिक समर्थन देण्याचा हा खरोखर आशादायक प्रयत्न वाटतो. पण याविषयी सांगण्यासारखे येथे मावेल त्याहून खूप अधिक आहे. तो सफल होतो की नाही, हे आता तुम्ही स्वतः ठरवावे.\footnote{पुढे वाचायचे असल्यास इनकुर्वाती (2020, पृ.~95--100) उपयुक्त ठरेल.}







\section{परिशिष्ट: प्रतिस्थापनाशी संबंधित निष्कर्ष}\label{sth:replacement:refproofs:sec}

या विभागात आपण $\ZF$ मध्ये प्रतिबिंबन सिद्ध करू. प्रतिबिंबन आणि प्रतिस्थापन कोणत्या अर्थाने समतुल्य आहेत हेही सिद्ध करू. तसेच प्रतिबिंबन/प्रतिस्थापन यांच्या सामर्थ्याविषयी या सर्वांचा एक रोचक परिणाम सिद्ध करू. \emph{सूचना: या पाठ्यपुस्तकातील गणिताचा हा सहजपणे सर्वांत प्रगत भाग आहे.}

सुरुवातीला एक साहाय्यक सिद्धांत पाहू. चलांच्या किंवा वस्तूंच्या अनुक्रमांसाठी संक्षेप म्हणून तो \emph{वरची रेघ} ही लेखनपद्धत वापरतो. म्हणजे ``$\overline{a}_k$'' हे ``$a_{k_1}$, \dots,
$a_{k_n}$'' यांचे संक्षिप्त लेखन आहे; येथे $n$ संदर्भावरून ठरतो.

\begin{lem}\label{sth:replacement:refproofs:lemreflection}
प्रत्येक $1 \leq i \leq k$ साठी $\phi_i(\overline{v}_i, x)$ हे सूत्र असू द्या. मग प्रत्येक $\alpha$ साठी काही $\beta > \alpha$ असा असतो की, कोणतेही $\overline{a}_1,
\ldots, \overline{a}_k \in V_\beta$ आणि प्रत्येक $1 \leq i \leq k$ यांच्यासाठी:
\[
  \exists x\phi_i(\overline{a}_i, x) \rightarrow (\exists x \in V_\beta) \phi_i(\overline{a}_i, x)
\]
\end{lem}

\begin{proof}
पुढीलप्रमाणे $\mu$ हे पद परिभाषित करू: $\mu(\overline{a}_1, \ldots,
\overline{a}_k)$ हा पुढील सर्व अटी, प्रत्येक $1 \leq i \leq k$ साठी, पूर्ण करणारा लघुतम टप्पा $V$ आहे:
\[
\exists x\phi_i(\overline{a}_i, x) \rightarrow (\exists x \in V) \phi_i(\overline{a}_i, x)
\]
सर्व $\overline{a}_1, \ldots, \overline{a}_k$ साठी $\mu(\overline{a}_1, \ldots, \overline{a}_k)$ अस्तित्वात आहे, हे सहज तपासता येते. आता प्रतिस्थापन आणि आपले पुनरावर्तन प्रमेय वापरून परिभाषित करू:
\begin{align*}
  S_0 & = V_{\alpha+1}\\
  S_{n+1} & = S_n \cup \bigcup
  \Setabs{\mu(\overline{a}_1, \ldots, \overline{a}_k)}
  {\overline{a}_1, \ldots, \overline{a}_k \in S_n} \\
  S &= \bigcup_{m < \omega} S_m.
\end{align*}
प्रत्येक $S_n$, आणि त्यामुळे $S$ सुद्धा, $V_\alpha$ नंतरचा टप्पा आहे. आता $\overline{a}_1$, \dots,~$\overline{a}_k \in S$ निश्चित करा; म्हणून काही $n <
\omega$ साठी $\overline{a}_1$, \dots, $\overline{a}_k \in S_n$. काही $1 \leq i \leq k$ निश्चित करा आणि $\exists x
\phi_i(\overline{a}_i,x)$ मानू. मग रचनेनुसार $(\exists x \in \mu(\overline{a}_1,
\ldots, \overline{a}_k))\phi_i(\overline{a}_i, x)$; म्हणून $(\exists x \in S_{n+1})\phi_i(\overline{a}_i, x)$ आणि परिणामी $(\exists x \in S)\phi_i(\overline{a}_i, x)$. म्हणून हाच $S$ आपला $V_\beta$ आहे.
\end{proof}
\noindent
आता \ref{sth:replacement:ref:reflectionschema} सहज सिद्ध करता येतो:

\begin{proof}[पुरावा]
$\alpha$ निश्चित करा. व्यापकता न गमावता, $\phi$ मध्ये फक्त $\exists$, $\lnot$ आणि $\land$ हेच तार्किक संयोगक आहेत असे मानता येते (कारण अभिव्यक्तीसाठी ते पुरेसे आहेत). $\psi_1, \ldots, \psi_k$ या $\phi$ च्या सर्व उपसूत्रांची जटिलतेनुसार यादी असू द्या; त्यामुळे $\psi_k = \phi$. \ref{sth:replacement:refproofs:lemreflection} नुसार असा $\beta > \alpha$ आहे की, कोणतेही $\overline{a}_i \in V_\beta$ आणि प्रत्येक $1 \leq i \leq k$ साठी:
\begin{align}\label{reflectionnicelybehaved}
  \exists x\psi_i(\overline{a}_i, x) \rightarrow
  (\exists x \in V_\beta) \psi_i(\overline{a}_i, x)\tag{*}
\end{align}
$\psi_i$ च्या जटिलतेवर विगमन करून, कोणत्याही $\overline{a}_i \in V_\beta$ साठी $\psi_i(\overline{a}_i) \leftrightarrow
\psi_i^{V_\beta}(\overline{a}_i)$ हे दाखवू. $\psi_i$ अणुसूत्र असेल, तर हे उघड आहे. ही द्विशर्त असेही दाखवते की, हा गुणधर्म असलेल्या उपसूत्रांचे निषेध किंवा संधीकरण $\psi_i$ असेल, तर $\psi_i$ लाही हा गुणधर्म असतो. म्हणून संख्यकाचे प्रकरणच लक्ष देण्यासारखे आहे. $\overline{a}_i \in V_\beta$ निश्चित करा; मग:
\begin{align*}
  (\exists x \psi_i(\overline{a}_i, x))^{V_\beta}
  &\text{ तेव्हा आणि केवळ तेव्हाच }
  (\exists x \in V_\beta)\psi_i^{V_\beta}(\overline{a}_i, x)
  &&\text{व्याख्येनुसार}\\
  &\text{ तेव्हा आणि केवळ तेव्हाच }
  (\exists x \in V_\beta)\psi_i(\overline{a}_i,  x)
  &&\text{विगमनगृहीतकानुसार}\\
  &\text{ तेव्हा आणि केवळ तेव्हाच }
  \exists x \psi_i(\overline{a}_i, x)
  &&\text{\eqref{reflectionnicelybehaved} नुसार}
\end{align*}
याने विगमन पूर्ण होते; $\psi_k = \phi$ असल्याने अपेक्षित निष्कर्ष मिळतो.
\end{proof}

आपण $\ZF$ मध्ये प्रतिबिंबन सिद्ध केले आहे. आपला पुरावा मूलतः मॉन्टेग्यू (1961) यांच्या पुराव्यानुसार आहे. आता $\Z$ मध्ये प्रतिबिंबनावरून प्रतिस्थापन मिळते हे सिद्ध करायचे आहे. हा पुरावा लेव्ही (1960) यांच्या पुराव्यावर आधारित आहे, पण थोडा सोपा केला आहे.

आपण $\Z$ मध्ये काम करत असल्याने वर दिलेल्या अचूक रूपात प्रतिबिंबन मांडता येणार नाही. कारण प्रतिबिंबन मांडताना आपण ``$V_\alpha$'' हे लेखन वापरले, आणि त्याची व्याख्या $\Z$ मध्ये करता येत नाही (\ref{sth:spine:zf:sec} पाहा). म्हणून त्याऐवजी प्रतिबिंबनाचे वरकरणी दुर्बल रूप पुढीलप्रमाणे मांडू:

\begin{defish}
\emph{दुर्बल प्रतिबिंबन.} कोणतेही सूत्र $\phi$ घ्या. मग असा संक्रमक संच $S$ असतो की $0$, $1$ आणि $\phi$ मधील कोणतीही प्राचले हे $S$ चे सदस्य असतात, आणि $(\forall \overline{x} \in S)(\phi \liff
\phi^S)$.
\end{defish}

यावरून प्रतिस्थापन सिद्ध करण्यासाठी आधी लेव्ही (1960, प्रमेय 2 चा पहिला भाग) यांचे अनुसरण करून, एकाच वेळी दोन सूत्रांचे ``प्रतिबिंबन'' करता येते हे दाखवू:

\begin{lem}[$\Z + \text{दुर्बल प्रतिबिंबन}$ मध्ये.]\label{sth:replacement:refproofs:lem:reflect}
कोणतीही सूत्रे $\psi, \chi$ घ्या. मग असा संक्रमक संच $S$ असतो की $0$ आणि $1$ (आणि सूत्रांमधील कोणतीही प्राचले) हे $S$ चे सदस्य असतात, आणि $(\forall \overline{x} \in S)((\psi \liff \psi^S) \land (\chi
\liff \chi^S))$.
\end{lem}

\begin{proof}
$\phi$ हे $(z = 0 \land \psi) \lor (z = 1 \land \chi)$ सूत्र असू द्या.

येथे आपण संक्षिप्त लेखन वापरले आहे: ``$z = 0$'' चे ``$\forall t\, t \notin z$'' असे, आणि ``$z =1$'' चे ``$\forall s(s \in z
\liff \forall t\, t \notin s)$'' असे पूर्ण लेखन करायला हवे. पण $0, 1 \in S$ आणि $S$ संक्रमक असल्याने, ही सूत्रे $S$ साठी \emph{निरपेक्ष} आहेत; म्हणजे त्यांचे संख्यक $S$ पुरते मर्यादित केले तरी ती त्याच वस्तूला लागू होतात.\footnote{अधिक औपचारिकपणे, $\xi$ हे या दोन सूत्रांपैकी कोणतेही एक असेल, तर $\xi(z) \liff \xi^S(z)$.}

दुर्बल प्रतिबिंबनानुसार योग्य असा काही $S$ आपल्याला मिळतो की:
\begin{align*}
  (\forall z, \overline{x} \in S)(&\phi \liff \phi^S)\\
  \text{म्हणजे }(\forall z, \overline{x} \in S)(&((z = 0 \land \psi) \lor (z = 1 \land \chi)) \liff {}\\
  &\phantom{(}((z = 0 \land \psi) \lor (z = 1 \land \chi))^S)\\
  \text{म्हणजे }(\forall z, \overline{x} \in S)(&((z = 0 \land \psi) \lor (z = 1 \land \chi))\liff {}\\
  &\phantom{(}((z = 0 \land \psi^S) \lor (z = 1 \land \chi^S)))\\
  \text{म्हणजे }(\forall \overline{x} \in S)(&(\psi \liff \psi^S) \land (\chi \liff \chi^S))
\end{align*}
दुसऱ्या विधानावरून तिसरे मिळते; कारण ``$z = 0$'' आणि ``$z=1$'' ही $S$ साठी निरपेक्ष आहेत. $0 \neq 1$ असल्याने चौथे विधान मिळते.
\end{proof}\noindent आता लेव्ही (1960, प्रमेय 6) यांचे अनुसरण करून, आणि त्यांचा पुरावा थोडा सोपा करून, प्रतिस्थापन मिळवू:

\begin{thm}[$\Z$ + दुर्बल प्रतिबिंबन मध्ये]\label{thm:replacement}
कोणतेही सूत्र $\phi(v,w)$ आणि कोणताही $A$ घ्या. $(\forall x \in A)\lexists![y][\phi(x,y)]$ असेल, तर $\Setabs{y}{(\exists x \in A)\phi(x,y)}$ अस्तित्वात असतो.
\end{thm}

\begin{proof}
$(\forall x \in A)\lexists![y][\phi(x,y)]$ पूर्ण करणारा $A$ निश्चित करा आणि सूत्रे परिभाषित करा:
\begin{align*}
  \psi &\text{ म्हणजे } (\phi(x, z) \land A = A)\\
  \chi &\text{ म्हणजे } \lexists[y][\phi(x, y)]
\end{align*}
$A$ हे $\psi$ चे प्राचल असल्याने, \ref{sth:replacement:refproofs:lem:reflect} वापरून असा संक्रमक~$S$ मिळतो की $0, 1, A \in S$ (तसेच इतर कोणतीही प्राचले त्या संचात असतात) आणि:
\[
  (\forall x,z \in S)((\psi \liff \psi^S) \land (\chi \liff \chi^S))
\]
म्हणून विशेषतः:
\begin{align*}
  (\forall  x, z \in S)(&\phi(x, z) \liff \phi^S(x, z))\\
  (\forall x \in S)(&\exists y\phi(x, y) \liff (\exists y \in S)\phi^S(x, y))
\end{align*}
हे एकत्र करून, आणि $A \in S$ व $S$ संक्रमक असल्यामुळे $A \subseteq S$ हे लक्षात घेऊन:
\begin{align*}
  (\forall x \in A)(&\exists y\phi(x, y) \liff (\exists y \in S)\phi(x, y))
\end{align*}
आता $(\forall x \in A)(\lexists![y \in S])\phi(x, y)$; कारण $(\forall x \in A)\lexists![y][\phi(x, y)]$. म्हणून विभक्तीकरणावरून $\Setabs{y \in S}{(\exists x \in A) \phi(x, y)} = \Setabs{y}{(\exists
x \in A) \phi(x, y)}$ मिळतो.
\end{proof}








\section{परिशिष्ट: सांत स्वयंसिद्धकीकरणीयता}\label{sth:replacement:finiteaxiomatizability:sec}
प्रतिस्थापनापासून काही निष्कर्ष काढून हे प्रकरण संपवू. पहिला निष्कर्ष मॉन्टेग्यू (1961) यांचा आहे. लक्षात घ्या: हा $\ZF$ च्या \emph{अंतर्गत} पुरावा नाही, तर $\ZF$ \emph{विषयीचा} पुरावा आहे:
\begin{thm}\label{sth:replacement:finiteaxiomatizability:zfnotfinitely}
  $\ZF$ सुसंगत असेल, तर ती सांत स्वयंसिद्धकांनी मांडता येत नाही. अधिक सर्वसाधारणपणे: $\Th{T}$ सांत असेल आणि $\Th{T} \Proves \ZF$ असेल, तर $\Th{T}$ विसंगत असते.

  (येथे आपण गृहीत धरतो की विचाराधीन प्रथम-क्रम वाक्यांमध्ये $\in$ हेच एकमेव अतार्किक मूलचिन्ह आहे; तसेच प्रत्येक $\phi \in \ZF$ साठी $\Th{T} \Proves \phi$ आहे हे दाखवण्यासाठी $\Th{T} \Proves \ZF$ असे लिहितो.)
\end{thm}
\begin{proof}
  $\Th{T} \Proves \ZF$ पूर्ण करणारी सांत $\Th{T}$ निश्चित करा. म्हणून $\Th{T}$ प्रतिबिंबन, म्हणजे \ref{sth:replacement:ref:reflectionschema}, सिद्ध करते. $\Th{T}$ सांत असल्याने तिचे एका संधीकरणाच्या रूपात, $\theta$ असे, पुनर्लेखन करता येते. $\Th{T}$ अनंतता सिद्ध करत असल्याने, अनंततेचे स्वयंसिद्धकही $\theta$ मध्ये संधीघटक म्हणून घेऊ; त्यामुळे $\theta^X$ पूर्ण करणारा $X$ रिक्त नसतो. या सूत्रावर प्रतिबिंबन वापरल्यास $\Th{T} \Proves \exists \beta(\theta \liff \theta^{V_\beta})$. स्पष्टपणे $\Th{T} \Proves \theta$ असल्याने $\Th{T} \Proves \exists \beta\ \theta^{V_\beta}$ मिळते.

  आता $\psi(X)$ हे पुढीलचे संक्षिप्त लेखन असू द्या:
  \[
    \theta^X \land X\text{ संक्रमक आहे} \land (\forall Y \in X)(Y\text{ संक्रमक आहे}\lif \lnot \theta^{Y})
  \]
साधारणपणे याचा अर्थ: $X$ हे $\theta$ चे संक्रमक प्रतिमान आहे आणि या बाबतीत $\in$-लघुतम आहे. आता $\Th{T} \Proves \exists \beta\ \theta^{V_\beta}$ हे आठवा. $\ZF$ मधील, आणि त्यामुळे $\Th{T}$ मधील, प्रतांकांविषयीच्या मूलभूत तथ्यांवरून:
  \begin{equation}
    \Th{T} \Proves \exists M \psi(M). \tag{*}\label{Mpsi}
  \end{equation}
$\psi(X)$ चे पहिले संधीघटक वापरल्यास, $\Th{T} \Proves \sigma$ असेल तेव्हा $\Th{T} \Proves \forall X(\psi(X) \lif \sigma^X)$ मिळते. म्हणून \eqref{Mpsi} वरून:
  \begin{align*}
    \Th{T} &\Proves \forall X(\psi(X) \lif (\exists N \psi(N))^X)\\
  \intertext{हे आणि पुन्हा \eqref{Mpsi} वापरल्यास:}
    \Th{T} &\Proves \exists M(\psi(M) \land (\exists N \psi(N))^M)
  \intertext{म्हणून विशेषतः:}
    \Th{T} &\Proves \exists M(\psi(M) \land (\exists N \in M)((N\text{ संक्रमक आहे})^M \land (\theta^N)^M))
  \intertext{आता संक्रमणाविषयीच्या प्राथमिक युक्तिवादावरून:}
    \Th{T} &\Proves \exists M(\psi(M) \land (\exists N \in M)(N\text{ संक्रमक आहे} \land \theta^N))
  \end{align*}
म्हणून $\Th{T}$ विसंगत आहे.\footnote{या ``प्राथमिक युक्तिवादा''त संक्रमक संचांसाठी काही ``निरपेक्षतेविषयीची तथ्ये'' सिद्ध करावी लागतात.}
\end{proof}

पॉटर (2004, पृ.~223) यांनी नोंदवलेला पुढील निष्कर्ष यासारखाच आहे:

\begin{prop}\label{sth:replacement:finiteaxiomatizability:finiteextensionofZ}
  $\Th{T}$ ही $\Z$ मध्ये सांत नवीन स्वयंसिद्धके जोडून मिळालेली उपपत्ती असू द्या. $\Th{T} \Proves \ZF$ असेल, तर $\Th{T}$ विसंगत आहे. (येथे \ref{sth:replacement:finiteaxiomatizability:zfnotfinitely} प्रमाणेच गृहीत निर्बंध आहेत.)
\end{prop}
\begin{proof}
  विभक्तीकरणाची (अनंत अनेक) उदाहरणे \emph{वगळून} $\Th{T}$ च्या सर्व स्वयंसिद्धकांचे संधीकरण $\theta$ ने दर्शवा. या पुराव्यासाठी $\psi_*(X)$ हे पुढीलचे संक्षिप्त लेखन असू द्या:
\[
  \theta^X \land X\text{ संक्रमक व उपसंच-समावेशक आहे}
  \land (\forall Y \in X)(Y\text{ संक्रमक व उपसंच-समावेशक आहे}\lif\lnot\theta^Y).
\]
उपसंच-समावेशकत्वाची व्याख्या \ref{sth:spine:Valphabasic:sec} मध्ये दिली आहे. प्रतिबिंबनातून मिळणारे टप्पे या दोन्ही अटी पूर्ण करतात. म्हणून \ref{sth:replacement:finiteaxiomatizability:zfnotfinitely} च्या प्रतांक-युक्तिवादात या अटी वापरून $\Th{T}\Proves\exists M\psi_*(M)$ मिळते.

  \ref{sth:replacement:finiteaxiomatizability:zfnotfinitely} प्रमाणे, $\Th{T} \Proves \sigma$ असेल तेव्हा $\Th{T} \Proves \forall X(\psi_*(X) \lif \sigma^X)$ मिळते हा रूपबंध सिद्ध करता येतो. मग \ref{sth:replacement:finiteaxiomatizability:zfnotfinitely} प्रमाणेच पुरावा पूर्ण होतो; येथे संक्रमक व उपसंच-समावेशक $M$ मध्ये आतील $N$ साठी या दोन्ही अटींची निरपेक्षता वापरतो.

  तथापि, हा रूपबंध सिद्ध करण्यासाठी \ref{sth:replacement:finiteaxiomatizability:zfnotfinitely} पेक्षा थोडे अधिक काम लागते. कारण विभक्तीकरणाची उदाहरणे $\Th{T}$ मध्ये असली, तरी ती $\theta$ ची संधीघटके नाहीत. पण प्रत्येक विभक्तीकरण-उदाहरण $\sigma$ साठी $\Th{T} \Proves \forall X(X\text{ संक्रमक व उपसंच-समावेशक आहे} \lif \sigma^X)$ हे सिद्ध करून हा अडथळा दूर करता येतो. हे वाचकांसाठी सोडतो.
\end{proof}
\begin{prob}
  प्रत्येक विभक्तीकरण-उदाहरण $\sigma$ साठी $\Z \Proves \forall X(X\text{ संक्रमक व उपसंच-समावेशक आहे} \lif \sigma^X)$ आहे हे दाखवा. (हा रूपबंध आपण \ref{sth:replacement:finiteaxiomatizability:finiteextensionofZ} मध्ये वापरला.)
\end{prob}
\begin{prob}
  प्रत्येक $\phi \in \Z$ साठी $\ZF \Proves \phi^{V_{\omega+\omega}}$ आहे हे दाखवा.
\end{prob}
\begin{prob}
  \ref{sth:replacement:finiteaxiomatizability:zfnotfinitely} आणि \ref{sth:replacement:finiteaxiomatizability:finiteextensionofZ} यांच्या पुराव्यांत वापरलेले उर्वरित रूपबंधात्मक निष्कर्ष तपासा.
\end{prob}

\ref{sth:replacement:absinf:sec} मध्ये म्हटल्याप्रमाणे, यावरून प्रतिस्थापन हे \ref{sth:ordinals:ordtype:thmOrdinalRepresentation} पेक्षा काटेकोरपणे अधिक सामर्थ्यवान आहे असे दिसते. किंवा थोडे अधिक नेमके सांगायचे तर: $\Z$ + ``प्रत्येक सुक्रमण एका अद्वितीय क्रमसूचक संख्येशी समरूप आहे'' ही उपपत्ती सुसंगत असेल, तर प्रतिस्थापनाचे काही उदाहरण तिला सिद्ध करता येत नाही.














